% Generated draft: requires mathematical and source review.
\documentclass[11pt]{article}
\usepackage[margin=25mm]{geometry}
\usepackage{fontspec}
\setmainfont{Latin Modern Roman}
\usepackage{amsmath,amssymb,mathtools,unicode-math}
\setmathfont{Latin Modern Math}
\usepackage{enumitem,needspace,xurl,hyperref}
\setlength{\parindent}{0pt}
\setlength{\parskip}{5pt}
\setlength{\emergencystretch}{3em}
\newcommand{\N}{\mathbb{N}}
\newcommand{\Z}{\mathbb{Z}}
\newcommand{\Q}{\mathbb{Q}}
\newcommand{\R}{\mathbb{R}}
\newcommand{\C}{\mathbb{C}}
\newcommand{\F}{\mathbb{F}}
\newcommand{\A}{\mathbb{A}}
\newcommand{\PP}{\mathbb{P}}
\newcommand{\E}{\mathbb{E}}
\newcommand{\eps}{\varepsilon}
\newcommand{\dd}{\,\mathrm{d}}
\begin{document}
\section*{UnsolvedMath ID: 701510}
\textbf{Source problem code:} catalogue entry
\subsection{Definitions and mathematical statement}
Let \(p\) be a prime. A \(p\)-group means an abstract group in which every element has finite order equal to a power of \(p\); the group itself need not be finite. A group \(G\) is finitely generated if \(G=\langle S\rangle\) for some finite subset \(S\subseteq G\). It is residually finite if, for every \(x\in G\setminus\{1\}\), there exist a finite group \(F\) and a homomorphism \(\varphi:G\to F\) such that \(\varphi(x)\ne 1\). A maximal subgroup is a proper subgroup \(M<G\) for which there is no subgroup \(H\) with \(M<H<G\).

The question is whether, for every prime \(p\), every finitely generated residually finite \(p\)-group \(G\), and every maximal subgroup \(M<G\), one necessarily has
\[
 gMg^{-1}=M\qquad\text{for every }g\in G.
\]
Thus the question concerns normality of all maximal proper subgroups in the abstract subgroup lattice; no topology or finite-index restriction on \(M\) is stipulated.
\subsection{Short English statement}
Must every maximal subgroup of a finitely generated residually finite \(p\)-group be normal? The group may be infinite, and \(p\) ranges over all primes.
\subsection{Sources}
\begin{itemize}
\item Supplied local UnsolvedMath registry record: permanent ID \(701510\), exact source problem code \texttt{catalogue entry}, title \emph{Kourovka Notebook Problem 18.81}; upstream revision \texttt{d371ee917ac9647fbb06cd6fc529deec6f3dda82}.
\item The supplied record attributes Problem 18.81 to D. S. Passman, proposed in 2014, and cites \emph{The Kourovka Notebook (Unsolved Problems in Group Theory)}, 18th issue (2014), with the 21st issue edition (2026), identified as \texttt{arXiv:1401.0300v48}. The collection metadata names E. I. Khukhro and V. D. Mazurov as editors.
\item The registry labels the problem \texttt{open}. This is supplied source metadata, not an independently verified assessment of its current mathematical status. The cited edition and attribution have not been independently checked for this draft.
\end{itemize}
\end{document}
